\section{Reproducibility and Computation Record}
%% ============================================================

\subsection{Provenance of the Reported Outputs}

The results archive contains complete final PPO checkpoints, job configurations, learning curves, probe summaries, and suite metadata. The archived Python source hashes and all 20 final checkpoint hashes were checked against the manifest and job files when preparing the main text. The executable binary itself is not distributed in this results archive; its recorded hash alone does not establish that a newly compiled binary is byte-identical.

The recorded environment uses Python 3.13.15, NumPy 2.1.3, PyTorch 2.11.0+cu128, Gymnasium 1.3.0, and device \texttt{cuda}. These are archived versions, not a new compatibility guarantee.

\begin{table}[t]
\centering
\caption{Archived cryptographic hashes for reproducibility verification.}
\label{tab:hashes}
\small
\begin{tabular}{@{}ll@{}}
\toprule
Component & SHA256 \\
\midrule
Posterior draws & \texttt{a2cd4f268b3edf0d...202d1b95} \\
Native extension & \texttt{8c1fe82a4830d09e...ef882780} \\
PPO training source & \texttt{12ffcf180b6fc899...662d3f15} \\
\bottomrule
\end{tabular}
\end{table}

The diagnostic gate records maximum rank-normalized split $\hat{R} = 1.008945$ and minimum bulk/tail ESS $= 606.86$, satisfying the software convergence gates of $1.01$ and $400$ \citep{vehtari2021rank}. Passing these thresholds addresses MCMC sampling quality, not Gaussianity of real-world monitoring errors or external market calibration. This supplement does not reconstruct missing likelihood inputs or invent posterior moments from those two statistics. Independent calibration replication requires the original labeled observations, observation-scale settings, sampling configuration, and retained draws matching the recorded posterior hash.

\subsection{Interaction Accounting}

\begin{table}[t]
\centering
\caption{Recorded interaction counts across all completed jobs.}
\label{tab:interactions}
\small
\begin{tabular}{@{}lrr@{}}
\toprule
Component & Per Job & Across 20 Jobs \\
\midrule
Training              & \num{1000000}  & \num{20000000} \\
Training diagnostics  & \num{819200}   & \num{16384000} \\
Trainer final eval.   & \num{8192}     & \num{163840} \\
Suite final eval.     & \num{32768}    & \num{655360} \\
Paired deviation audit & \num{1048576} & \num{20971520} \\
\midrule
Total recorded        & \num{2908736}  & \num{58174720} \\
\bottomrule
\end{tabular}
\end{table}

Each audit type costs \num{131072} transitions, including the pre-intervention prefix in both branches. Eight types give the reported \num{1048576}. One market transition contains both firms' actions; it is not two separate environment interactions. Recorded completed-job times sum to $21.978$ hours for baseline and $22.079$ hours for ablation ($44.06$ hours total). Hardware-specific runtime comparisons require additional telemetry beyond these job totals.

Three previous-attempt records concern baseline seed 20261010 and ablated seeds 20261008 and 20261009. All are marked running, not failed, in the archived records. Their filenames contain \texttt{previous\_failure}, but the records do not establish a numerical training failure or its cause. All three final jobs completed, and unrecorded interrupted work is excluded from the interaction and runtime totals above. Finally, the broader research plan also scheduled Soft Actor-Critic (SAC) conditions, but no incomplete SAC comparison enters this study; because the PPO-only scope and its interpretation were finalized after inspecting completed runs, the manuscript is not a preregistered confirmatory test of a null effect.

%% ============================================================
